\chapter*{译读附记}
\addcontentsline{toc}{chapter}{译读附记}
以下说明补充记录第三篇中沿用的原书记号及表述。页码均指原书，右侧链接返回相应正文。第一、二篇的实质性校订及疑点已在相关页的译注中说明；第一篇的正式勘误另见前一附录。

\begin{description}
\item[第 180 页，延拓 $F$ 的叙述]
上文取 $F=K_C/(2g-2)$，这里延拓时直接写 $\omega_{X/B}$；译文沿用这一记号，未另添归一化因子。
\hfill\hyperref[orig:185]{返回正文}
\item[第 181 页，引理 8.1]
陈述中未显式假设 $X$ 紧致，证明中则使用了紧致性。本次补充了测度非退化条件的说明，见\hyperlink{translator-note-22}{译者补注22}。\hfill\hyperref[orig:186]{返回正文}
\item[第 182 页，引理 8.2]
第二项的分母为 $2d_1$，紧接的文字则给出 $(2g-1)/r$；两处均沿用原书。\hfill\hyperref[orig:187]{返回正文}
\item[第 183 页，正合列附近]
$a,b$ 既用于整除关系 $ba^2=d$，也用于直和重数；译文沿用原书记号。\hfill\hyperref[orig:188]{返回正文}
\item[第 187 页，推论 11.4]
原书定义 $\mathscr P_{ij}=(\operatorname{pr}_i+\operatorname{pr}_j)^*L-\operatorname{pr}_i^*L-\operatorname{pr}_j^*L$，后面对 $(s_i\operatorname{pr}_i-s_j\operatorname{pr}_j)^*L$ 的展开末项仍印为 $+s_is_j\mathscr P_{ij}$；译文保留该加号。\hfill\hyperref[orig:192]{返回正文}
\item[第 189 页，引理 12.1(a)]
文字称常数 $c_1,c_2$，公式右端却使用 $c_3^{\delta}$；译文保留原书的对应字母。\hfill\hyperref[orig:194]{返回正文}
\item[第 193 页，式 (13.11)]
指数中 $c\sum d_i h(X_i)+c''\sum d_i$ 的加号沿用原书。\hfill\hyperref[orig:198]{返回正文}
\item[第 196 页，Hensel 引理的证明]
原书 Taylor 展开的一个式子写为 $f(\alpha')=f(\alpha)-J_f(\alpha)(\alpha'-\alpha)+\cdots$，一次项前的减号予以保留。\hfill\hyperref[orig:201]{返回正文}
\item[第 197 页，引理 15.8]
原文称 $J_f$ 为 Jacobian determinant，但定义和后文使用的是前述矩阵；译文写作“如引理 15.3 中的 Jacobian”。\hfill\hyperref[orig:202]{返回正文}
\item[第 198 页，归纳步骤 (i)]
求和下限原印为 $j=1$，后文 (15.11) 的对应求和则从 $0$ 开始；两处沿用原书。\hfill\hyperref[orig:203]{返回正文}
\item[第 200 页，参考文献引用]
正文写 [So 1]，文末相应文献项标为 [So]；译文保留原有引用键。\hfill\hyperref[orig:205]{返回正文}
\item[第 201 页，引理 16.2 的证明]
原书一处列 $f_1,\ldots,f_{m_i}$，前文则为 $f_1,\ldots,f_{M-m_i}$；译文保留两处原有范围。\hfill\hyperref[orig:206]{返回正文}
\item[第 204 页，第 17 节]
原书所写 $\deg\bigl((s_i^2\operatorname{pr}_i-s_j^2\operatorname{pr}_j)^*L\bigr)|_E$ 的态射系数含平方，紧接右端则含 $s_i^2,s_j^2$；译文保留这一原式。\hfill\hyperref[orig:209]{返回正文}
\end{description}
